\clearpage
\begin{trapbox}
\begin{center}\begin{tikzpicture}[font=\small,>=stealth]
\draw[gray!45,rounded corners=2pt] (0,0) rectangle (14.2,8.5);
\node at (7.1,7.85) {同样 $y_0=0$，先分清中心与非中心};
\node at (7.1,7) {两图同为 $x^2/4+y^2=1$，比例相同。};
\foreach \j/\cx in {1/3.5,2/10.5}{
 \begin{scope}[shift={(\cx,4.4)},scale=1.1]
 \draw[gray,->] (-2.35,0)--(2.35,0) node[right] {$x$};
 \draw[blue!75!black,thick] (0,0) ellipse(2 and 1);
 \ifnum\j=1
  \draw[orange!85!black,thick] (-2,0)--(2,0);
  \draw[purple!80!black,thick] (-1.897366596101,-.316227766017)--(1.897366596101,.316227766017);
  \draw[gray,dashed,thick] (0,-1)--(0,1);
  \fill (0,0) circle(1.8pt) node[below right=6pt] {$O=M$};
 \else
  \draw[orange!85!black,very thick] (1,-.8660254038)--(1,.8660254038);
  \fill (1,.8660254038) circle(1.6pt) node[above right=5pt] {$A$};
  \fill (1,-.8660254038) circle(1.6pt) node[below right=5pt] {$B$};
  \fill (1,0) circle(1.8pt) node[below left=5pt] {$M(1,0)$};
 \fi
 \end{scope}}
\node at (3.5,2) {$M=O$：水平、竖直、倾斜都可以};
\node at (10.5,2) {$M\ne O$：此例只能取 $x=1$};
\node at (7.1,.75) {把 $0/0$ 当成“竖直”，会丢掉左图中的其他方向。};
\end{tikzpicture}
\end{center}
\begin{description}[style=nextline,leftmargin=2.8em,labelsep=0.5em,itemsep=0.45em]
\item[\textbf{\textcolor{CTrap}{陷阱一：一看到 $y_0=0$ 就说弦竖直}}]
\textbf{正确做法：}先分清 $M=O$ 还是 $M\ne O$。
对椭圆：$M=O$ 时任何过中心的直径都可以；$M=(x_0,0)$、$0<|x_0|<a$ 时非退化弦是 $x=x_0$。
对双曲线：$M=O$ 时可取过中心的割线（例如 $y=0$）；$M=(x_0,0)$、$|x_0|>a$ 时可取竖直弦 $x=x_0$。
对抛物线：$M=(x_0,0)$、$x_0>0$ 时非退化弦为竖直弦 $x=x_0$。
\textbf{最小反例：}椭圆中 $M=O$，水平弦 $y=0$ 存在，所以“$y_0=0$ 必竖直”错误。

\end{description}
\tcbbreak
\begin{center}% Same hyperbola and isotropic scale .4 in both panels.
\begin{tikzpicture}[font=\small,>=stealth]
\draw[gray!45,rounded corners=2pt] (0,0) rectangle (14.2,10.5);
\node at (7.1,9.85) {先有候选方向，再看直线能否真正截出两点};
\node at (7.1,8.95) {两图同为 $x^2/4-y^2/5=1$，$Q(x,y)=x^2/4-y^2/5$，横纵同尺度。};
\foreach \j/\cx in {1/3.3,2/10.7}{
 \begin{scope}[shift={(\cx,4)},scale=.4]
 \draw[gray,->] (-6.5,0)--(7,0) node[right] {$x$};
 \draw[gray,->] (0,-4.2)--(0,7.1) node[above] {$y$};
 \draw[blue!75!black,thick,domain=-4:7,samples=100,variable=\v]
  plot ({2*sqrt(1+\v*\v/5)},\v);
 \draw[blue!75!black,thick,domain=-4:7,samples=100,variable=\v]
  plot ({-2*sqrt(1+\v*\v/5)},\v);
 \ifnum\j=1
  \pgfmathsetmacro{\u}{6*sqrt(330)/55}
  \coordinate (A) at ({4-\u},{3-5*\u/3});
  \coordinate (B) at ({4+\u},{3+5*\u/3});
  \draw[orange!85!black,very thick] (A)--(B);
  \fill[orange!85!black] (A) circle(3pt);
  \fill[orange!85!black] (B) circle(3pt);
  \fill[orange!85!black] (4,3) circle(3.2pt) node[right=6pt] {$M_1$};
  \draw[purple!80!black,very thick] ({-sqrt(24/5)},1)--({sqrt(24/5)},1);
  \fill[purple!80!black] ({-sqrt(24/5)},1) circle(3pt);
  \fill[purple!80!black] ({sqrt(24/5)},1) circle(3pt);
  \fill[purple!80!black] (0,1) circle(3.2pt) node[above left=6pt] {$M_2$};
 \else
  \draw[orange!85!black,dashed,thick] (-2.8,-3.75)--(5.8,7);
  \fill (1,1) circle(3.2pt) node[above=10pt] {$N$};
 \fi
 \end{scope}}
\node[text=orange!85!black] at (3.3,2) {$M_1(4,3),\quad Q(M_1)=11/5>1$};
\node[text=purple!80!black] at (3.3,1.2) {$M_2(0,1),\quad Q(M_2)=-1/5<0$};
\node[text=teal!80!black] at (3.3,.5) {两种位置都能产生实弦};
\node at (10.7,2) {$N(1,1),\quad Q(N)=1/20\in(0,1)$};
\node[text=orange!85!black] at (10.7,1.2) {候选线：$y-1=\tfrac54(x-1)$};
\node at (10.7,.5) {没有实交点，更谈不上中点弦};
\end{tikzpicture}
\end{center}
\begin{description}[style=nextline,leftmargin=2.8em,labelsep=0.5em,itemsep=0.45em]
\item[\textbf{\textcolor{CTrap}{陷阱二：写出斜率就默认两个交点存在}}]
\textbf{正确做法：}点差法是\emph{必要条件}，反求弦要补查“严格不等号”：
\[
\begin{array}{ll}
\text{椭圆：}&\dfrac{x_0^2}{a^2}+\dfrac{y_0^2}{b^2}<1;\\[4pt]
\text{抛物线：}&2px_0-y_0^2>0;\\[4pt]
\text{双曲线（$y_0\ne0$）：}&Q<0\ \text{或}\ Q>1,\quad
Q=\dfrac{x_0^2}{a^2}-\dfrac{y_0^2}{b^2}.
\end{array}
\]
\textbf{最小反例：}双曲线 $x^2/4-y^2/5=1$、$M(1,1)$ 算得 $k=5/4$，但 $Q=1/20$，根本不存在这样的实弦。

\item[\textbf{\textcolor{CTrap}{陷阱三：椭圆、双曲线的符号记反}}]
\textbf{正确做法：}不背“椭负双正”也行。两式相减后
$\dfrac{x_0}{a^2}\Delta x\ \boldsymbol{\pm}\ \dfrac{y_0}{b^2}\Delta y=0$；
原方程是 $+y^2$，移项得负号；原方程是 $-y^2$，移项得正号。

\item[\textbf{\textcolor{CTrap}{陷阱四：抛物线的 $p$ 与一次项系数混淆}}]
\textbf{正确做法：}题目若是 $y^2=4x$，则 $2p=4\Rightarrow p=2$，斜率为 $k=2/y_0$。
$y^2=2px$ 下焦点是 $(p/2,0)$；一般写成 $y^2=4fx$，就有 $p=2f$，两种记号绝不可混用。

\end{description}
\tcbbreak
\begin{center}\begin{tikzpicture}[font=\small,>=stealth]
\draw[gray!45,rounded corners=2pt] (0,0) rectangle (14.2,9);
\node at (7.1,8.35) {圆的垂直性质，要在伸缩后的坐标中使用};
\node at (7.1,7.45) {$X=x/2,\quad Y=y/\sqrt3$：椭圆变单位圆，中点仍对应中点。};
\pgfmathsetmacro{\u}{sqrt(20/21)}
\foreach \j/\cx in {1/3.5,2/10.5}{
 \begin{scope}[shift={(\cx,4.1)}]
 \ifnum\j=1
  \coordinate (A) at ({1-\u},{1+.75*\u});
  \coordinate (B) at ({1+\u},{1-.75*\u});
  \coordinate (M) at (1,1);
  \draw[blue!75!black,thick] (0,0) ellipse (2 and 1.7320508076);
  \def\xl{x}\def\yl{y}
 \else
  % Display scale 2cm per X/Y unit; actual transformed coordinates are x/2,y/sqrt3.
  \coordinate (A) at ({1-\u},{2*(1+.75*\u)/sqrt(3)});
  \coordinate (B) at ({1+\u},{2*(1-.75*\u)/sqrt(3)});
  \coordinate (M) at (1,{2/sqrt(3)});
  \draw[blue!75!black,thick] (0,0) circle(2);
  \def\xl{X}\def\yl{Y}
 \fi
 \draw[gray,->] (-2.4,0)--(2.45,0) node[right] {$\xl$};
 \draw[gray,->] (0,-2.15)--(0,2.3) node[above] {$\yl$};
 \draw[gray] (2,-.05)--(2,.05);
 \ifnum\j=1
  \node[below right=4pt] at (2,0) {$2$};
 \else
  \node[below right=4pt] at (2,0) {$1$};
 \fi
 \draw[orange!85!black,very thick] (A)--(B);
 \draw[purple!80!black,thick] (0,0)--(M);
 \fill (0,0) circle(1.5pt) node[below left=4pt] {$O$};
 \ifnum\j=1
  \fill (A) circle(1.6pt) node[above left=5pt] {$A$};
  \fill (B) circle(1.6pt) node[above right=5pt] {$B$};
  \fill (M) circle(1.8pt) node[below=8pt] {$M$};
 \else
  \draw[gray] (1.151186,1.023770)--(1.020255,.872584)--(.869069,1.003514);
  \fill (A) circle(1.6pt) node[above left=5pt] {$A'$};
  \fill (B) circle(1.6pt) node[above right=5pt] {$B'$};
  \fill (M) circle(1.8pt) node[below=8pt] {$M'$};
 \fi
 \end{scope}}
\node at (3.5,1.4) {原坐标：$k_{OM}=1,\ k_{AB}=-3/4$};
\node at (10.5,1.4) {伸缩后：$OM'\perp A'B'$};
\node at (7.1,.55) {非等比伸缩会改变角度，不能把圆中的垂直关系直接搬回椭圆。};
\end{tikzpicture}
\end{center}
\begin{description}[style=nextline,leftmargin=2.8em,labelsep=0.5em,itemsep=0.45em]
\item[\textbf{\textcolor{CTrap}{陷阱五：把“圆的垂直”生搬到椭圆}}]
\textbf{错误想法：}认为原坐标里 $OM\perp AB$，直接写 $k=-x_0/y_0$。
\textbf{正确做法：}椭圆需先按 $X=x/a,\ Y=y/b$ 伸缩成圆，再使用垂直性质；
原坐标应有系数 $b^2/a^2$。
如 $x^2/4+y^2/3=1, M(1,1)$，正确 $k=-3/4$，并不是 $-1$。

\end{description}
\tcbbreak
\begin{center}\begin{tikzpicture}[font=\small,>=stealth]
\draw[gray!45,rounded corners=2pt] (0,0) rectangle (14.2,11.4);
\node at (7.1,10.75) {中点满足 $y_0=1$，但只落在一条开射线上};
\node at (7.1,9.9) {在 $y^2=4x$ 上画三条斜率均为 $2$ 的实弦。};
\begin{scope}[shift={(3.2,4)},scale=1]
\draw[gray,->] (-.5,0)--(5.3,0) node[right] {$x$};
\draw[gray,->] (0,-2.2)--(0,4.4) node[above] {$y$};
\draw[blue!75!black,thick,domain=-2.1:4.2,samples=100,variable=\v]
 plot ({\v*\v/4},\v);
\foreach \m/\col in {.75/orange!85!black,1.5/purple!80!black,2.5/teal!80!black}{
 \pgfmathsetmacro{\v}{sqrt(4*\m-1)}
 \draw[\col,thick] ({\m-\v/2},{1-\v})--({\m+\v/2},{1+\v});
 \fill[\col] ({\m-\v/2},{1-\v}) circle(1.6pt);
 \fill[\col] ({\m+\v/2},{1+\v}) circle(1.6pt);
 \fill[purple!80!black] (\m,1) circle(2.3pt);
}
\draw[purple!80!black,very thick,->] (.25,1)--(5.2,1);
\filldraw[fill=white,draw=purple!80!black,thick] (.25,1) circle(2.5pt)
 node[left=8pt,text=black] {$T$};
\end{scope}
\node[text=blue!75!black] at (10.9,8.2) {$y^2=4x$};
\node[text=purple!80!black] at (10.9,6.8) {中点在 $y=1$ 上移动};
\node at (10.9,5.6) {$T=(1/4,1)$ 为空心点};
\node at (10.9,4.5) {到 $T$ 时两端点重合};
\node at (10.9,3.2) {三条弦都满足 $k=2$};
\node at (7.1,1.15) {$\displaystyle \boxed{y_0=1,\quad x_0>\frac14}\qquad$切点边界不取，左侧也不可达。};
\end{tikzpicture}
\end{center}
\begin{description}[style=nextline,leftmargin=2.8em,labelsep=0.5em,itemsep=0.45em]
\item[\textbf{\textcolor{CTrap}{陷阱六：把“中点轨迹方程”当作完整轨迹}}]
\textbf{正确做法：}关系式只是轨迹所满足的方程，通常还须去掉切点边界和不可达部分。
如 $y^2=4x$ 的斜率 $k=2$ 的中点轨迹，是 $y=1,\ x>1/4$，\textbf{不是}整条 $y=1$。
\end{description}
\end{trapbox}
\clearpage
